\subsection{DEFINITION OF AN ORDER RELATION}

Let \(\mathbb{R}\{x, y\}\) be a relation, \(x\) and \(y\) being distinct letters. \(\mathbb{R}\) is said to be an order relation with respect to the letters \(x\) and \(y\) (or between \(x\) and \(y\) ) if

\[
(R \{x, y \} \text {   and   } R \{y, z \}) \Longrightarrow R \{x, z \},
\]

\[
(\mathrm{R} \{x, y \} \text {   and   } \mathrm{R} \{y, x \}) \Longrightarrow (x = y),
\]

\[
\mathrm{R} \left\{x, y \right\} \Longrightarrow (\mathrm{R} \left\{x, x \right\} \text { and } \mathrm{R} \left\{y, y \right\}).
\]

The first of the above relations says that R is transitive with respect to the letters x and y (Chapter II, §6, no. 1).

Examples

\begin{enumerate}
\def\labelenumi{(\arabic{enumi})}
\item
  The relation of equality, \(x = y\) , is an order relation.
\item
  The relation \(X \subset Y\) is an order relation between X and Y (Chapter II, §1, no. 2, Propositions 1 and 2 and axiom A1) which is often called the inclusion relation.
\item
  Let \(\mathbf{R}\{x, y\}\) be an order relation between \(x\) and \(y\) . The relation \(\mathbf{R}\{y, x\}\) is then an order relation between \(x\) and \(y\) , called the opposite of the order relation \(\mathbf{R}\{x, y\}\) .
\end{enumerate}

An order relation on a set E is an order relation \(R\{x, y\}\) with respect to two distinct letters x, y such that the relation \(R\{x, x\}\) is equivalent to \(x \in E\) (in other words, is such that \(R\{x, y\}\) is reflexive on E (Chapter II, §6, no. 1)). Then the relation \(R\{x, y\}\) implies `` \(x \in E\) and \(y \in E\) '' and the relation \((R\{x, y\}\) and \(R\{y, x\})\) is equivalent to `` \(x \in E\) and \(y \in E\) and \(x = y\) ''.

Examples

\begin{enumerate}
\def\labelenumi{(\arabic{enumi})}
\item
  The relations of equality and inclusion are not order relations on a set, for the relations \(x = x\) and \(\mathbf{X} \subset \mathbf{X}\) are not collectivizing (Chapter II, § 1, no. 7).
\item
  Let \(\mathbb{R}\{x, y\}\) be an order relation between \(x\) and \(y\) , and let \(\mathbb{E}\) be a set such that \(x \in \mathbb{E}\) implies \(\mathbb{R}\{x, x\}\) (notice that the empty set satisfies this condition). The relation `` \(\mathbb{R}\{x, y\}\) and \(x \in \mathbb{E}\) and \(y \in \mathbb{E}\) '' is then an order relation on \(\mathbb{E}\) , as is immediately verified; it is called the order relation induced by \(\mathbb{R}\{x, y\}\) on \(\mathbb{E}\) (cf.~no. 4). By abuse of language, the phrase ``the relation \(S\{x, y\}\) is an order relation between elements of \(\mathbb{E}\) '' is often used in place of ``the relation ( \(S\{x, y\}\) and \(x \in \mathbb{E}\) and \(y \in \mathbb{E}\) ) is an order relation on \(\mathbb{E}\) ''. For example, if \(A\) is a set, the relation `` \(X \subset Y\) and \(X \subset A\) and \(Y \subset A\) '' is an order relation between subsets of \(A\) .
\item
  Let E, F be sets. The relation ``g extends f'' is an order relation (between f and g) on the set of mappings of subsets of E into F.
\item
  In the set \(\mathfrak{P}(\mathfrak{P}(E))\) of sets of subsets of a set \(E\) , let \(\mathfrak{F}\) be the set of partitions of \(E\) (Chapter II, § 4, no. 7). We recall that a partition \(\varpi\) is said to be coarser than a partition \(\varpi'\) if given any \(Y \in \varpi'\) there exists \(X \in \varpi\) such that \(Y \subset X\) (Chapter II, § 4, no. 6). For each partition \(\varpi\) of \(E\) let \(\tilde{\varpi}\) be the graph of the equivalence relation defined by \(\varpi\) on \(E\) (Chapter II, § 6, no. 2), that is to say, the union of the (mutually disjoint) sets \(A \times A\) , where \(A\) runs through \(\varpi\) . The relation `` \(\varpi\) is coarser than \(\varpi'\)'' is immediately seen to be equivalent to \(\tilde{\varpi} \supset \tilde{\varpi}'\) , and is therefore an order relation on the set \(\mathfrak{F}\) between \(\varpi\) and \(\varpi'\) .
\end{enumerate}

An ordering on a set E is a correspondence \(\Gamma = (G, E, E)\) with E as source and as target, and such that the relation \((x, y) \in G\) is an order relation on E. By abuse of language we shall sometimes refer to the graph G of \(\Gamma\) as an ordering on E. If \(R\{x, y\}\) is an order relation on E, it has a graph which is an ordering on E.

PROPOSITION 1. A correspondence \(\Gamma\) between E and E is an ordering on E if and only if its graph G satisfies the following conditions:

\begin{enumerate}
\def\labelenumi{(\alph{enumi})}
\item
  \(\mathbf{G} \circ \mathbf{G} = \mathbf{G}\) .
\item
  The set \(\mathbf{G} \cap \overline{\mathbf{G}}\) is the diagonal \(\Delta\) of \(\mathbf{E} \times \mathbf{E}\) .
\end{enumerate}

For the relation \(((x, y) \in \mathbf{G}\) and \((y, z) \in \mathbf{G}) \Longrightarrow ((x, z) \in \mathbf{G})\) can be written as \(\mathbf{G} \circ \mathbf{G} \subset \mathbf{G}\) , and the relation

\((x, y) \in G\) and \((y, x) \in G) \iff (x = y\) and \(x \in E\) and \(y \in E)\)

can be written as \(\mathbf{G} \cap \overline{\mathbf{G}} = \Delta\) . From \(\mathbf{G} \cap \overline{\mathbf{G}} = \Delta\) we then deduce \(\Delta \subset \mathbf{G}\) , whence \(\mathbf{G} = \Delta \circ \mathbf{G} \subset \mathbf{G} \circ \mathbf{G}\) ; since also \(\mathbf{G} \circ \mathbf{G} \subset \mathbf{G}\) , we have

\[
\mathbf {G} \circ \mathbf {G} = \mathbf {G}.
\]

